\hypertarget{chapter-17-title-pages-front-matter-and-print-ready-output}{%
\chapter{Title Pages, Front Matter, and Print-Ready Output}\label{chapter-17-title-pages-front-matter-and-print-ready-output}}

A title page is the first thing any reader sees, and it is also, for exactly that reason, one of the last things an author actually finalizes. This chapter covers the design and technical mechanics of professional front matter and print-ready output, but it opens with a tension worth naming directly, because it shapes how the rest of the chapter's advice should be applied rather than being a side issue to mention and move past.

\hypertarget{a-genuine-contradiction-early-metadata-late-titles}{%
\section{A Genuine Contradiction: Early Metadata, Late Titles}\label{a-genuine-contradiction-early-metadata-late-titles}}

LaTeX's structure asks an author to declare a document's title, author, and date near the top of the source, in the preamble, before a single word of the actual content has been written, using \texttt{\textbackslash{}title}, \texttt{\textbackslash{}author}, and \texttt{\textbackslash{}date}, or the class-level equivalents from the previous chapter. This placement is not arbitrary; the preamble is also where packages are loaded and where the document's typographic identity, fonts, theorem styles, notation, is fixed, and a class expecting required front-matter fields, as the previous chapter recommended, naturally wants those fields declared early, in the same place as everything else the document depends on structurally.

The difficulty is that a book's actual title is often one of the last things to stabilize, not the first, particularly for a long-form work whose central argument develops and sometimes shifts meaningfully over the course of drafting. A working title chosen at the outset frequently turns out, by the time a manuscript is finished, to describe an earlier and narrower version of the book than what was actually written, and retitling at that point is not merely a matter of changing the string inside \texttt{\textbackslash{}title\{\}}. A title change of any real consequence tends to reflect, or force, a change in emphasis that ought to be felt throughout the book: which chapter gets foregrounded in the preface, which thread gets called the central argument in the introduction and which gets recast as supporting material, occasionally even which appendix material deserves promotion into the main text now that the book's own sense of what it is about has shifted.

There is no purely technical resolution to this, because the difficulty is not really about where \texttt{\textbackslash{}title\{\}} sits in the file. It is a genuine mismatch between LaTeX's front-loaded declaration of a document's identity and the fact that a book's identity, in the sense that actually matters, is frequently not settled until close to the end of writing it. The practical response is to treat the title declared early as provisional by default for any project expected to take more than a few weeks, rather than as a commitment made in the preamble and then forgotten about; and, more importantly, to treat a late title change as a prompt to reread the introduction and preface specifically for passages that were written to support the old title's framing rather than the new one, since those are exactly the passages most likely to have drifted out of alignment with what the book, retitled, is now claiming to be about. A title change that touches only the preamble and nothing else in the manuscript is a signal worth treating with some suspicion: it may mean the new title is genuinely a surface relabeling with no deeper implications, but it may also mean the rest of the book has not yet caught up to what the new title is promising a reader.

\hypertarget{designing-a-title-page}{%
\section{Designing a Title Page}\label{designing-a-title-page}}

With that caution in mind, the mechanics of a title page itself are straightforward once a custom \texttt{\textbackslash{}makebooktitle}, as described in the previous chapter, is in place. A title page that reads as professionally published rather than templated typically uses restrained typography, the same type family as the body text rather than a decorative display face reserved only for this one page, generous whitespace rather than a page crowded with every possible piece of metadata, and a clear typographic hierarchy: the title set larger and with more weight than the subtitle, the subtitle larger than the author line, the author line larger than the affiliation. Centering everything by default, without considering whether a left-aligned or asymmetric layout might better suit the specific book, is one of the most common tells of an untailored title page; the choice should be made deliberately, matching the visual conventions established by the rest of the book's design, rather than defaulted to.

\hypertarget{front-matter-and-back-matter-conventions}{%
\section{Front Matter and Back Matter Conventions}\label{front-matter-and-back-matter-conventions}}

Front matter, the material preceding the main text proper, conventionally includes a half-title page, the full title page, a copyright and edition page, a dedication where relevant, a table of contents, and, where present, lists of figures and tables generated automatically from the document's own float captions rather than maintained by hand. A preface, distinct from an introduction, addresses the reader about the book itself, its motivation and scope, while an introduction begins the book's actual argument; conflating the two, or omitting the preface entirely for a book that would benefit from directly addressing why it exists before diving into its subject matter, is a common structural gap in self-published technical work that a deliberate front-matter plan avoids.

Back matter, following the main text, conventionally includes appendices, a bibliography, an index, a glossary or nomenclature list where the project maintains one as described in Chapter 13, and, for a book meant to read as a complete, finished object, a colophon: a short closing note on how the book was produced, the typeface used, the engine it was built with. Placed correctly, a colophon belongs entirely in the back matter, never on the title page itself; a stray production note or build instruction accidentally left on a title page is not merely an aesthetic slip; automated tools that ingest and summarize documents have been observed treating such a note as though it were meaningful content of the book rather than production metadata, which is reason enough to keep this material strictly confined to its proper place in the back matter.

\hypertarget{cover-design-and-print-ready-pdfs}{%
\section{Cover Design and Print-Ready PDFs}\label{cover-design-and-print-ready-pdfs}}

A PDF intended only for screen reading has different requirements than one intended for print, and conflating the two produces a file that looks fine on a monitor and fails, sometimes expensively, at a print shop. Print-ready output requires bleed, extra content extending past the trim line on any page with full-bleed color or imagery, so that minor variation in the physical cutting process does not leave a visible white sliver at the edge; a defined trim size matching the exact physical dimensions the printer expects, distinct from the page size a screen-reading PDF would use; and a color profile appropriate to the print process, typically CMYK for offset printing rather than the RGB most digital tools default to, since colors that look correct in RGB on a screen can shift unpredictably when converted to CMYK ink without deliberate profile management.

Font embedding is the other requirement that print workflows enforce far more strictly than screen-reading ones: every font used anywhere in the document, including any custom or licensed font from Chapter 10, has to be fully embedded in the final PDF, not merely referenced by name, since a print shop's rendering pipeline has no access to fonts installed only on the author's own machine. LuaLaTeX and XeLaTeX both embed fonts automatically as part of ordinary PDF generation, but confirming this with a PDF inspection tool before sending a file to print, rather than assuming it happened correctly, catches the rare case where a font's own licensing metadata restricts embedding in a way that silently produces a PDF referencing a font that will not actually be present at the print shop's end.

Cover design itself is typically produced separately from the interior typesetting, in a dedicated design tool rather than LaTeX, since a wraparound cover with spine width dependent on final page count is a fundamentally different kind of layout problem than interior typesetting; LaTeX can compute an accurate final page count for exactly this purpose, \texttt{\textbackslash{}pageref\{LastPage\}} from the \texttt{lastpage} package, which is often the one piece of information a cover designer actually needs from the interior's own build process before finalizing spine width.

With title pages, front matter, and print-ready output covered, the next chapter turns to publishing pipelines proper: producing more than one output format, a book, a standalone paper excerpt, a website, from a single maintained source tree.
